\documentclass[10pt]{article}

\usepackage[utf8]{inputenc}

\usepackage[T1]{fontenc}

\usepackage{amsmath}

\usepackage{amsfonts}

\usepackage{amssymb}

\usepackage[version=4]{mhchem}

\usepackage{stmaryrd}

\usepackage{bbold}

\usepackage{graphicx}

\usepackage[export]{adjustbox}

\graphicspath{ {./images/} }



\begin{document}

видным после перехода к нормализуюшим координатам $y$. Тор имеет вид прямого произведения $\boldsymbol{r}$ эллипсов, лежащих на 2-мерных попарно ортогональных плоскостях $\mathbb{E}_{y_{i}}^{2}, \dot{y}_{i}$ Эти эллипсы - уровни энергий



\[

E_{i}=\left(\dot{y}_{i}^{2}+\omega_{i}^{2} y_{i}^{2}\right) / 2=\omega_{i}^{2}\left|\xi^{i}\right|^{2} D_{i}^{2} / 2

\]



составляющих собственных колебаний.



Траектории точек $x=\varphi(t)$ в конфигурационном пространстве $\mathbb{E}_{x}^{n}$ называются фи гурам и Лиссажу. Они являются проекциями на $\mathbb{E}_{x}^{n}$ обмоток $r$-мерных торов из $\mathbb{E}_{x, p}^{2 n}$. Траектория собственного колебания является проекцией эллипса и имеет вид отрезка. При несоизмеримых частотах $\omega_{i}=\sqrt{\lambda_{i}}$ на торе (то есть если $(k, \omega)=\sum k_{i} \omega_{i}=0$ для всех $\left.k \in \mathbb{Z}^{r} \backslash\{0\}\right)$ траектория $x=\varphi(t)$ заметает множество, замыкание к-рого есть $r$-мерный параллелепипед. Его стороны параллельны натянутым на векторы $\xi^{i}$ главным осям квадратичной формы $U=(A x, x) / 2$. Полное его заполнение связано с тем, что в нерезонансном случае каждая обмотка тора всюду плотна на нем.



Прямое обобщение уравнения $\ddot{x}=-A x$ линейных колебаний - гамильтоново линейное уравнение, задаваемое диагонализуемым оператором, спектр к-рого лежит на мнимой оси (так наз. эллиптич. случай, см. [6], [1]). (О нелинейных малых собственных колебаниях см., напр., [5].) Для исследования нелинейных М.к. с произвольными начальными условиями используют приведение их в окрестности особой точки при негамильтоновом возмущении к нормальным формам Пуанкаре - Дюлака, а при гамильтоновом - к нормальным формам Биркгофа (см. [3], [1], [6], [4]). При этом приведение, как правило, формальное, то есть на уровне, вообще говоря, расходящихся рядов. Однако и это позволяет получать определенную информацию об устойчивости колебаний.



Чрезвычайно полезным оказался подход ко многим линейным задачам математич. физики, как к линейным колебаниям с бесконечным числом степеней свободы. Простейший пример - смешанная задача Коши для волнового уравнения. Линейные колебания лежат в основе и квантовой теории поля.



Лит.: [1] Ар н о л ь д В. И., Математические методы классической механики, М., 1989; [2] е го же, Обыкновенные дифференциальные уравнения, М., 1984; [3] Арнольд В.И., Козлов В. В., Ней штадт А.И., в кн.: Итоги науки и техники, Современные проблемы математики. Фундаментальные направления, т. 3, М., 1985; [4] Карапетян А. В., Румянце в В. В., в кн.: Итоги науки и техники, Общая механика, т. 6, М., 1983; [5] Манев ич Л.И., М и хлин Ю. В., П ил и пчук В. Н., Метод нормальных колебаний для существенно нелинейных систем, М., 1989; [6] М о з р Ю., Лекции о гамильтоновых системах, пер. с англ., М., 1973.



H. H. Hexopoues.



МАЛЫЙ КАНОНИЧЕСКИЙ АНСАМБЛЬ - СМ. СТатистический ансамбль, Статистической механики математические задачи.



МАНДЕЛСТАМА ПРЕДСТАВЛЕНИЕ, Д В ОЙ Н О е спектральное представление,- простейшее интегральное представление для амплитуды рассеяния элементарных частиц (см. Дисперсионных соотношений метод) как функции инвариантных квадрата полной энергии $s$ в системе центра масс и квадрата передачи 4-импульса $t$.



В простейшем случае бинарного процесса $\mathrm{a}+\mathrm{b} \rightarrow \mathrm{c}+\mathrm{d}$ (рис. 1) упругого рассеяния частиц с равными массами $m$ (напр., двух пионов) М. п. имеет вид:



\[

\begin{gathered}

A(s, t)=\frac{1}{\pi^{2}} \iint d s^{\prime} d t^{\prime} \frac{\rho_{s t}\left(s^{\prime}, t^{\prime}\right)}{\left(s^{\prime}-s\right)\left(t^{\prime}-t\right)}+ \\

+\frac{1}{\pi^{2}} \iint d u^{\prime} d t^{\prime} \frac{\rho_{u t}\left(u^{\prime}, t^{\prime}\right)}{\left(u^{\prime}-u\right)\left(t^{\prime}-t\right)}+\frac{1}{\pi^{2}} \iint d u^{\prime} d s^{\prime} \frac{\rho_{s u}\left(s^{\prime}, u^{\prime}\right)}{\left(u^{\prime}-u\right)\left(s^{\prime}-s\right)},

\end{gathered}

\]



где релятивистски-инвариантные переменные



\[

s=\left(p_{\mathrm{a}}+p_{\mathrm{b}}\right)^{2}, t=\left(p_{\mathrm{a}}-p_{\mathrm{c}}\right)^{2}, u=\left(p_{\mathrm{a}}-p_{\mathrm{d}}\right)^{2}

\]



связаны друг с другом соотношением



\[

s+u+t=4 m^{2}

\]



( $p_{\mathrm{a}, \mathrm{b}, \mathrm{c}, \mathrm{d}}-4$-импульсы частиц $\mathrm{a}, \mathrm{b}, \mathrm{c}, \mathrm{d}$; используется система единиц, в к-рой $c=1$ ), а спектральные плотности $\rho$ отличны от нуля только в областях I, II, III (рис. 2), так что амплитуда $A$ аналитична при всех комплексных $s$ и $t$ за исключением этих вещественных областей. М. П. задает и аналитич. свойства амплитуды как функции одной комплексной переменной ( $s$ или $t$ ) - это разрезы, определяемые асимптотами границы спектральных функций; $s>4 m^{2}$, $t>4 m^{2}, u>4 m^{2}$.



Важным свойством



М. п. является его явная перекрестная симметрия: оно определяет единую аналитич. функщию, к-рая в разных областях переменных $s, u$ описывает различные перекрестные процессы



\includegraphics[max width=\textwidth, center]{2023_07_08_79956ea9c7b3a32ce80ag-1}

(рис. 2)



Представление предложено С. Манделстамом (S. Mandelstam) в 1959 и строго доказано в квантовой механике с потенциалом взаимодействия определенного класса. Характерной особенностью М. п. в этом случае является нулевое значение спектральной плотности $\rho_{s u}$. Однако в квантовой теории поля его удалось доказать лишь в рамках перенормированной теории возмущений.



\begin{center}

\includegraphics[max width=\textwidth]{2023_07_08_79956ea9c7b3a32ce80ag-1(1)}

\end{center}



Pис. 2. Области аналитичности амплитуды процесса $\mathrm{a}+\mathrm{b} \rightarrow \mathrm{c}+\mathrm{d}$ (в случае одинаковых масс частиц). I, II, III - области, где отличны от нуля спектральные плотности р. Заштрихованы области перекрестных процессов.



М. п., наряду с унитарности условием, составляет основу дисперсионного подхода в теории элементарных частиц. Связывая амплитуды различных процессов, оно приводит к системе нелинейных интегральных уравнений. Однако возникающая система оказывается настолько широкой, что включает в себя амплитуды практически всех процессов, происходящих с элементарными частицами, и не поддается математич. разрешению. В ряде случаев с помощью различных приближений удается сузить систему и получить интересные физич. результаты. М. п. прочно вошло в арсенал аналитич. методов теории элементарных частиц и лежит в основе многих моделей.



Лит.: [1] Ширков Д. В., Серебряков В. В., Мещеря к о в В. А., Дисперсионные теории сильных взаимодействий при низких энергиях, М., 1967; [2] Боголюбов Н. Н., Логуно в А.А., Тодо ров И. Т., Основы аксиоматического подхода в квантовой теории поля, М., 1969; [3] И и и ксо о К К., 3 ю бе е Ж Ж.-Б., Квантовая теория поля, пер. с англ., т. 1, М., 1984. А. В. Ефремов.





\end{document}